\documentclass[11pt]{article}
\usepackage[margin=1in]{geometry}
\usepackage{amsmath, amssymb, amsfonts}
\usepackage{enumitem}
\usepackage{hyperref}
\usepackage{xcolor}

\title{Comprehensive Calculus III Final Review \\ \large Problem Sets, Solutions, and Mental Roadmaps}
\author{Math Curriculum Generator}
\date{}

\begin{document}

\maketitle
\tableofcontents
\newpage

% ---------------------------------------------------------
% CONCEPT CHECKLIST & MAPPING
% ---------------------------------------------------------
\section{Concept Checklist \& Question Mapping}

Use this matrix to identify which questions target the specific competencies required for mastery of Calculus III concepts found in your review.

\begin{itemize}
    \item \textbf{Q1: Paraboloid Equations.} Finding equations based on vertex, orientation, and scaling.
    \item \textbf{Q2: Equations of Planes (Normal Vector).} Constructing planes from normals, parallel planes, and orthogonal intersections.
    \item \textbf{Q3: Equations of Planes (3 Points).} Using cross products of displacement vectors to find a normal vector.
    \item \textbf{Q4: Angles Between Planes.} Using the dot product of normal vectors.
    \item \textbf{Q5: Area of Parallelograms \& Triangles in 3D.} Utilizing the magnitude of the cross product.
    \item \textbf{Q6: Parametric Equations of Lines.} Creating lines from points, parallel vectors, and intersections.
    \item \textbf{Q7: First Partial Derivatives.} Applying product, quotient, and chain rules in multivariable contexts.
    \item \textbf{Q8: Second Partial Derivatives.} Computing $f_{xx}, f_{yy}$, and verifying Clairaut's Theorem ($f_{xy} = f_{yx}$).
    \item \textbf{Q9: Directional Derivatives.} Computing $D_{\mathbf{u}}f = \nabla f \cdot \mathbf{u}$, ensuring $\mathbf{u}$ is a unit vector.
    \item \textbf{Q10: Maximum Rate of Change (Gradient).} Utilizing the properties of the gradient vector for magnitude and direction.
    \item \textbf{Q11: Volume via Double Integrals (Cartesian).} Setting up and evaluating integrals over Type I/Type II regions.
    \item \textbf{Q12: Volume via Double Integrals (Polar).} Converting domains and integrands to $r$ and $\theta$, including the Jacobian $r$.
    \item \textbf{Q13: Cartesian to Cylindrical Conversion (Setup).} Translating $dx dy dz$ bounds into $dz dr d\theta$.
    \item \textbf{Q14: Cartesian to Cylindrical Conversion (Algebraic Integrands).} Handling complex Cartesian roots via $r$.
    \item \textbf{Q15: Evaluating Cylindrical Triple Integrals.} Finding volumes between paraboloids, planes, and spheres.
    \item \textbf{Q16: Cartesian to Spherical Conversion.} Translating complex integrands and bounds into $\rho, \theta, \phi$.
    \item \textbf{Q17: Solid Region Description in Spherical.} Understanding physical geometry (cones, spheres, xy-plane) in terms of angle boundaries.
    \item \textbf{Q18: Green's Theorem.} Converting closed loop line integrals into double integrals, paying strict attention to orientation.
\end{itemize}

\newpage
% ---------------------------------------------------------
% PROBLEM SETS
% ---------------------------------------------------------
\section{Problem Sets}

\subsection{Q1: Paraboloid Equations}
\begin{enumerate}[label=\textbf{1\alph*.}]
    \item \textbf{(Vanilla)} Find an equation of the paraboloid with the vertex at $(0, 2, -1)$, its axis parallel to the z-axis, and opening towards positive values of $z$. Assume the scaling factor $c = 1$.
    \item \textbf{(Twist 1 - Different Axis)} Find an equation of the paraboloid with the vertex at $(3, -4, 2)$, its axis parallel to the x-axis, opening towards negative values of $x$. Assume $c = 1$.
    \item \textbf{(Twist 2 - Passing through a point)} Find the equation of the paraboloid with vertex $(1, 1, 1)$, opening towards positive $y$ along an axis parallel to the y-axis, that passes through the point $(3, 5, 2)$.
    \item \textbf{(Harder Twist - Complete the Square)} Identify the vertex, the axis of symmetry, and the opening direction of the paraboloid given by the equation $z = -2x^2 + 12x - y^2 - 4y - 20$.
\end{enumerate}

\subsection{Q2: Equations of Planes (Normal Vector)}
\begin{enumerate}[label=\textbf{2\alph*.}]
    \item \textbf{(Vanilla)} Find the equation of the plane with normal vector $\mathbf{n} = \langle 3, -1, 4 \rangle$ passing through the point $(-2, 5, 1)$.
    \item \textbf{(Twist 1 - Parallel Plane)} Find the equation of the plane passing through the origin that is parallel to the plane $4x - 2y + 7z = 15$.
    \item \textbf{(Twist 2 - Orthogonal to a Line)} Find the equation of the plane passing through $(1, -1, 3)$ that is perpendicular to the line given by parametric equations $x = 2+3t$, $y = -t$, $z = 4+5t$.
    \item \textbf{(Harder Twist - Orthogonal to Two Planes)} Find the equation of the plane passing through $(2, 0, -3)$ that is perpendicular to both the planes $x + y - z = 2$ and $2x - y + 3z = 4$.
\end{enumerate}

\subsection{Q3: Equations of Planes (3 Points)}
\begin{enumerate}[label=\textbf{3\alph*.}]
    \item \textbf{(Vanilla)} Find an equation of the plane containing the points $P(2, 1, 1)$, $Q(-1, 3, 2)$, and $R(0, -2, 4)$.
    \item \textbf{(Twist 1 - Intercepts)} Find the equation of the plane whose x, y, and z intercepts are $3$, $-2$, and $5$ respectively.
    \item \textbf{(Twist 2 - Collinear Trap)} Determine if the points $A(1, 2, 3)$, $B(2, 4, 6)$, and $C(3, 6, 9)$ determine a unique plane. Why or why not?
    \item \textbf{(Harder Twist - Triangle Area)} For the plane containing $A(1, 0, 0)$, $B(0, 2, 0)$, and $C(0, 0, 3)$, find the plane's equation AND the area of the triangle $ABC$.
\end{enumerate}

\subsection{Q4: Angle Between Planes}
\begin{enumerate}[label=\textbf{4\alph*.}]
    \item \textbf{(Vanilla)} Find the angle between the planes $2x - y + z = 5$ and $x + y + 2z = 3$.
    \item \textbf{(Twist 1 - Coordinate Plane)} Find the angle between the plane $x + \sqrt{3}y + 2z = 4$ and the $xy$-plane.
    \item \textbf{(Twist 2 - Orthogonal Check)} Find the value of $k$ such that the planes $3x + ky - 2z = 7$ and $4x - y + kz = 1$ are orthogonal.
    \item \textbf{(Harder Twist - Parallel Check)} Find values for $a$ and $b$ such that the planes $ax + 4y - 6z = 2$ and $2x - by + 3z = 5$ are parallel.
\end{enumerate}

\subsection{Q5: Area of Parallelograms in 3D}
\begin{enumerate}[label=\textbf{5\alph*.}]
    \item \textbf{(Vanilla)} Find the area of the parallelogram spanned by the vectors $\mathbf{a} = \langle 2, -1, 4 \rangle$ and $\mathbf{b} = \langle 1, 3, -2 \rangle$.
    \item \textbf{(Twist 1 - From Points)} Find the area of the parallelogram with adjacent edges formed by the points $P(1, 1, 1)$, $Q(3, 2, 0)$, and $R(2, 4, 2)$ (where $P$ is the shared vertex).
    \item \textbf{(Twist 2 - Area of a Triangle)} Find the area of the triangle with vertices $(0, 0, 0)$, $(1, -2, 3)$, and $(-1, 4, 2)$.
    \item \textbf{(Harder Twist - Unknown Component)} Let $\mathbf{u} = \langle 1, 2, 0 \rangle$ and $\mathbf{v} = \langle 0, c, 3 \rangle$. If the area of the parallelogram they span is exactly $7$, find all possible values of $c$.
\end{enumerate}

\subsection{Q6: Parametric Equations of a Line}
\begin{enumerate}[label=\textbf{6\alph*.}]
    \item \textbf{(Vanilla)} Find the parametric equations of the line passing through the points $(2, 4, -1)$ and $(5, 0, 3)$.
    \item \textbf{(Twist 1 - Parallel Line)} Find the parametric equations of the line passing through $(1, 1, 1)$ that is parallel to the line $x = 3-t$, $y = 2t$, $z = -4+5t$.
    \item \textbf{(Twist 2 - Orthogonal to Plane)} Find the parametric equations of the line passing through $(-3, 2, 6)$ that is perpendicular to the plane $4x - y + 7z = 12$.
    \item \textbf{(Harder Twist - Line of Intersection)} Find the direction vector of the line formed by the intersection of the planes $x + y + z = 1$ and $x - 2y + 3z = 1$.
\end{enumerate}

\subsection{Q7: First Partial Derivatives}
\begin{enumerate}[label=\textbf{7\alph*.}]
    \item \textbf{(Vanilla)} Find $f_x$ and $f_y$ for $f(x,y) = x^4 y^3 - e^{2x} + \ln(y)$.
    \item \textbf{(Twist 1 - Product/Chain)} Find $f_x$ and $f_y$ for $f(x,y) = x^2 \sin(xy^3)$.
    \item \textbf{(Twist 2 - Quotients)} Find $f_x$ and $f_y$ for $f(x,y) = \frac{2x - y}{x + 3y}$.
    \item \textbf{(Harder Twist - 3 Variables)} Find $f_x, f_y, f_z$ for $f(x,y,z) = e^{xyz} \cos(x^2 + z)$.
\end{enumerate}

\subsection{Q8: Second Partial Derivatives}
\begin{enumerate}[label=\textbf{8\alph*.}]
    \item \textbf{(Vanilla)} Find all four second partial derivatives ($f_{xx}, f_{yy}, f_{xy}, f_{yx}$) for $f(x,y) = x^3 y^4 - 2xy$.
    \item \textbf{(Twist 1 - Exponential)} Find $f_{xy}$ for $f(x,y) = y e^{x^2y}$.
    \item \textbf{(Twist 2 - Clairaut's Theorem)} Verify that $f_{xy} = f_{yx}$ for $f(x,y) = \ln(x^2 + y^2)$.
    \item \textbf{(Harder Twist - Laplace Equation)} Show that $u(x,y) = e^x \sin y$ satisfies the Laplace equation $u_{xx} + u_{yy} = 0$.
\end{enumerate}

\subsection{Q9: Directional Derivatives}
\begin{enumerate}[label=\textbf{9\alph*.}]
    \item \textbf{(Vanilla)} Find the directional derivative of $f(x,y) = x^2y - 3y^2$ at the point $(1, 2)$ in the direction of the vector $\mathbf{v} = \langle 3, 4 \rangle$.
    \item \textbf{(Twist 1 - Given Angle)} Find the directional derivative of $f(x,y) = y e^{-x}$ at $(0, 4)$ in the direction given by the angle $\theta = 2\pi/3$.
    \item \textbf{(Twist 2 - Towards a Point)} Find the directional derivative of $f(x,y) = \sqrt{xy}$ at $(2, 8)$ in the direction pointing towards the origin $(0,0)$.
    \item \textbf{(Harder Twist - 3D Directional)} Find the directional derivative of $f(x,y,z) = x^2y + y^2z$ at $(1, 2, -1)$ in the direction of $\mathbf{v} = 2\mathbf{i} - \mathbf{j} + 2\mathbf{k}$.
\end{enumerate}

\subsection{Q10: Maximum Rate of Change (Gradient)}
\begin{enumerate}[label=\textbf{10\alph*.}]
    \item \textbf{(Vanilla)} Find the maximum rate of change of $f(x,y) = x^3 - y^3$ at the point $(2, -1)$ and the direction in which it occurs.
    \item \textbf{(Twist 1 - Square Root Function)} Find the maximum rate of change of $f(x,y) = \sqrt{x^2 + 2y}$ at $(3, 8)$.
    \item \textbf{(Twist 2 - Minimum Rate of Change)} In what direction does $f(x,y) = x e^{xy}$ decrease the fastest at the point $(1, 0)$? What is that minimum rate of change?
    \item \textbf{(Harder Twist - Zero Rate of Change)} For $f(x,y) = x^2 + 4y^2$ at the point $(2, 1)$, find a unit vector direction $\mathbf{u}$ such that the directional derivative $D_{\mathbf{u}}f = 0$.
\end{enumerate}

\subsection{Q11: Double Integral Volume (Cartesian)}
\begin{enumerate}[label=\textbf{11\alph*.}]
    \item \textbf{(Vanilla)} Find the volume under the surface $z = x^2 + y^2$ and above the rectangle $R = \{(x,y) \mid 0 \le x \le 2, 0 \le y \le 3\}$.
    \item \textbf{(Twist 1 - Type I Region)} Find the volume under $z = xy$ and above the region bounded by $y = x^2$ and $y = 4x$.
    \item \textbf{(Twist 2 - Type II Region)} Evaluate $\iint_D e^{x/y} dA$ where $D$ is bounded by $y=1$, $y=2$, $x=0$, and $x=y^2$.
    \item \textbf{(Harder Twist - Reverse Order of Integration)} Evaluate the integral by reversing the order of integration: $\int_0^4 \int_{\sqrt{x}}^2 \frac{1}{y^3 + 1} dy \, dx$.
\end{enumerate}

\subsection{Q12: Double Integral Volume (Polar)}
\begin{enumerate}[label=\textbf{12\alph*.}]
    \item \textbf{(Vanilla)} Evaluate $\iint_R (x^2 + y^2) dA$ where $R$ is the region bounded by the circle $x^2+y^2=9$ in the first quadrant.
    \item \textbf{(Twist 1 - Upper Half Plane)} Find the volume under $z = e^{x^2+y^2}$ and above the region bounded by the upper half of the circle $x^2+y^2=16$ and the x-axis.
    \item \textbf{(Twist 2 - Annulus)} Evaluate $\iint_R \sin(x^2+y^2) dA$ where $R$ is the region between the circles $x^2+y^2=1$ and $x^2+y^2=4$.
    \item \textbf{(Harder Twist - Shifted Circle)} Set up the polar double integral to find the area of the region bounded by the circle $x^2+y^2 = 2x$.
\end{enumerate}

\subsection{Q13: Cartesian to Cylindrical Conversion (Setup)}
\begin{enumerate}[label=\textbf{13\alph*.}]
    \item \textbf{(Vanilla)} Convert $\int_{-2}^2 \int_0^{\sqrt{4-x^2}} \int_{0}^{x^2+y^2} x^2 dz dy dx$ into an integral in cylindrical coordinates.
    \item \textbf{(Twist 1 - Different Variables)} Convert $\int_{0}^3 \int_{-\sqrt{9-y^2}}^{\sqrt{9-y^2}} \int_{0}^{\sqrt{x^2+y^2}} xyz \, dz dx dy$ into cylindrical coordinates.
    \item \textbf{(Twist 2 - Sphere Bound)} Convert $\int_{-1}^1 \int_{-\sqrt{1-x^2}}^{\sqrt{1-x^2}} \int_{-\sqrt{1-x^2-y^2}}^{\sqrt{1-x^2-y^2}} (x^2+y^2) dz dy dx$ into cylindrical coordinates.
    \item \textbf{(Harder Twist - Shifted Cylinder Region)} Convert the triple integral $\iiint_E z dV$ into cylindrical coordinates where $E$ is bounded by the cylinder $x^2+y^2=2y$, below by $z=0$, and above by $z=x^2+y^2$.
\end{enumerate}

\subsection{Q14: Cartesian to Cylindrical Conversion (Algebraic)}
\begin{enumerate}[label=\textbf{14\alph*.}]
    \item \textbf{(Vanilla)} Convert $\int_0^3 \int_0^{\sqrt{9-x^2}} \int_{0}^{x^2+y^2} \frac{z}{\sqrt{x^2+y^2}} dz dy dx$ into cylindrical coordinates.
    \item \textbf{(Twist 1 - Full Circle)} Convert $\int_{-2}^2 \int_{-\sqrt{4-x^2}}^{\sqrt{4-x^2}} \int_{x^2+y^2}^{4} \sqrt{x^2+y^2} dz dy dx$ into cylindrical coordinates.
    \item \textbf{(Twist 2 - Exponentials)} Convert $\int_0^1 \int_0^{\sqrt{1-y^2}} \int_0^{1} e^{x^2+y^2} dz dx dy$ into cylindrical coordinates.
    \item \textbf{(Harder Twist - Cone limits)} Convert $\int_{-1}^1 \int_{0}^{\sqrt{1-y^2}} \int_{\sqrt{x^2+y^2}}^{1} z dz dx dy$ into cylindrical coordinates.
\end{enumerate}

\subsection{Q15: Evaluating Cylindrical Triple Integrals}
\begin{enumerate}[label=\textbf{15\alph*.}]
    \item \textbf{(Vanilla)} Evaluate $\iiint_V z dV$, where $V$ is the solid bounded by $z = 0$ and $z = 9 - x^2 - y^2$.
    \item \textbf{(Twist 1 - Cone and Plane)} Evaluate the volume of the solid bounded above by $z = 4$ and below by the cone $z = \sqrt{x^2+y^2}$.
    \item \textbf{(Twist 2 - Between two paraboloids)} Evaluate $\iiint_V (x^2+y^2) dV$ where $V$ is bounded by $z = x^2+y^2$ and $z = 8 - x^2 - y^2$.
    \item \textbf{(Harder Twist - Center of Mass)} Set up the integrals in cylindrical coordinates to find the $z$-coordinate of the center of mass of the uniform solid bounded by $z = \sqrt{x^2+y^2}$ and $z = 4$.
\end{enumerate}

\subsection{Q16: Cartesian to Spherical Conversion}
\begin{enumerate}[label=\textbf{16\alph*.}]
    \item \textbf{(Vanilla)} Convert $\int_{-2}^2 \int_{-\sqrt{4-x^2}}^{\sqrt{4-x^2}} \int_{0}^{\sqrt{4-x^2-y^2}} z^2 \sqrt{x^2+y^2+z^2} dz dy dx$ to spherical coordinates.
    \item \textbf{(Twist 1 - Different octant)} Convert $\int_0^3 \int_0^{\sqrt{9-x^2}} \int_{0}^{\sqrt{9-x^2-y^2}} x dz dy dx$ to spherical coordinates.
    \item \textbf{(Twist 2 - Cone lower bound)} Convert $\int_{-1}^1 \int_{-\sqrt{1-x^2}}^{\sqrt{1-x^2}} \int_{\sqrt{x^2+y^2}}^{\sqrt{2-x^2-y^2}} \frac{1}{x^2+y^2+z^2} dz dy dx$ to spherical coordinates.
    \item \textbf{(Harder Twist - Offset plane)} Convert $\iiint_E z dV$ to spherical coordinates, where $E$ is the solid bounded above by the sphere $x^2+y^2+z^2=4$ and below by the plane $z=1$.
\end{enumerate}

\subsection{Q17: Solid Region Description in Spherical}
\begin{enumerate}[label=\textbf{17\alph*.}]
    \item \textbf{(Vanilla)} Set up the triple integral in spherical coordinates for the volume of the solid bounded above by the sphere $\rho = 5$ and below by the cone $\phi = \pi/4$.
    \item \textbf{(Twist 1 - Between two cones)} Describe the solid bounded by the spheres $\rho=1$ and $\rho=3$, and between the cones $\phi = \pi/6$ and $\phi = \pi/3$ in the first octant.
    \item \textbf{(Twist 2 - Ice Cream Cone volume)} Evaluate the volume of the solid bounded below by the $xy$-plane, on the sides by the sphere $\rho = 6$, and above by the cone $\phi = \pi/6$. (Wait, physically if it is above $\phi=\pi/6$ and below $\phi=\pi/2$, it's a wedge).
    \item \textbf{(Harder Twist - Shifted Sphere)} Describe the region enclosed by the sphere $x^2+y^2+z^2 = 2z$ in spherical coordinates.
\end{enumerate}

\subsection{Q18: Green's Theorem}
\begin{enumerate}[label=\textbf{18\alph*.}]
    \item \textbf{(Vanilla)} Evaluate $\oint_C (3y) dx + (5x) dy$ where $C$ is the boundary of the rectangle $[0,2] \times [0,1]$ traversed counter-clockwise.
    \item \textbf{(Twist 1 - Clockwise Trap)} Evaluate $\oint_C (y^2) dx + (3xy) dy$ where $C$ is the circle $x^2+y^2=9$ traversed \textit{clockwise}.
    \item \textbf{(Twist 2 - Nasty Integrand)} Evaluate $\oint_C (e^{x^2} - 2y) dx + (4x + \sin(y^3)) dy$ where $C$ is the triangle with vertices $(0,0), (4,0), (0,3)$ traversed CCW.
    \item \textbf{(Hardest - Area Application)} Use Green's Theorem to find the area of the region enclosed by the parametric curve $x(t) = \cos^3 t$, $y(t) = \sin^3 t$ for $0 \le t \le 2\pi$.
\end{enumerate}

\newpage
% ---------------------------------------------------------
% DETAILED SOLUTIONS
% ---------------------------------------------------------
\section{Detailed Solutions}

\subsection{Solutions to Q1}
\textbf{1a:} Vertex $(h,k,l) = (0,2,-1)$. Axis parallel to z-axis means equation is $z - l = c(x-h)^2 + c(y-k)^2$. Opens positive $z \implies c > 0$. Let $c=1$. \textbf{Answer: $z + 1 = x^2 + (y-2)^2$.}\\
\textbf{1b:} Vertex $(3, -4, 2)$. Axis parallel to x-axis $\implies$ linear term is $x$. Opens negative $x \implies$ negative constants. \textbf{Answer: $x - 3 = -(y+4)^2 - (z-2)^2$.}\\
\textbf{1c:} Vertex $(1,1,1)$, linear $y$. Equation: $y-1 = a(x-1)^2 + b(z-1)^2$. Passes through $(3,5,2)$: $5-1 = a(3-1)^2 + b(2-1)^2 \implies 4 = 4a + b$. If circular paraboloid, $a=b \implies 5a=4 \implies a=4/5$. Assuming $a=b=c$, \textbf{Answer: $y - 1 = c(x-1)^2 + c(z-1)^2$ where $c$ solves specific bounds.}\\
\textbf{1d:} $z = -2(x^2 - 6x + 9) - (y^2 + 4y + 4) - 20 + 18 + 4 = -2(x-3)^2 - (y+2)^2 + 2$. \textbf{Vertex $(3, -2, 2)$, axis parallel to z-axis, opens downward (negative z).}

\subsection{Solutions to Q2}
\textbf{2a:} $3(x - (-2)) - 1(y - 5) + 4(z - 1) = 0 \implies \mathbf{3x - y + 4z = -7}$.\\
\textbf{2b:} Parallel means same normal vector $\langle 4, -2, 7 \rangle$. Through origin: $4(x-0) - 2(y-0) + 7(z-0) = 0 \implies \mathbf{4x - y + 7z = 0}$.\\
\textbf{2c:} Line vector $\mathbf{v} = \langle 3, -1, 5 \rangle$ becomes the plane's normal $\mathbf{n}$. $3(x-1) - 1(y+1) + 5(z-3) = 0 \implies \mathbf{3x - y + 5z = 19}$.\\
\textbf{2d:} $\mathbf{n_1} = \langle 1, 1, -1 \rangle, \mathbf{n_2} = \langle 2, -1, 3 \rangle$. Plane normal $\mathbf{n} = \mathbf{n_1} \times \mathbf{n_2} = \langle (3-1), (-2-3), (-1-2) \rangle = \langle 2, -5, -3 \rangle$. Plane: $2(x-2) - 5(y-0) - 3(z+3) = 0 \implies \mathbf{2x - 5y - 3z = 13}$.

\subsection{Solutions to Q3}
\textbf{3a:} $\vec{PQ} = \langle -3, 2, 1 \rangle, \vec{PR} = \langle -2, -3, 3 \rangle$. $\mathbf{n} = \vec{PQ} \times \vec{PR} = \langle 9, 7, 13 \rangle$. Plane: $9(x-2)+7(y-1)+13(z-1)=0 \implies \mathbf{9x+7y+13z=38}$.\\
\textbf{3b:} Intercepts formula: $\frac{x}{a} + \frac{y}{b} + \frac{z}{c} = 1$. \textbf{Answer: $\frac{x}{3} - \frac{y}{2} + \frac{z}{5} = 1$ or $10x - 15y + 6z = 30$.}\\
\textbf{3c:} $\vec{AB} = \langle 1, 2, 3 \rangle, \vec{AC} = \langle 2, 4, 6 \rangle$. Since $\vec{AC} = 2\vec{AB}$, vectors are parallel. \textbf{Points are collinear, they do not determine a unique plane.}\\
\textbf{3d:} $\vec{AB} = \langle -1, 2, 0 \rangle, \vec{AC} = \langle -1, 0, 3 \rangle$. $\mathbf{n} = \langle 6, 3, 2 \rangle$. Plane: $\mathbf{6x+3y+2z=6}$. Area = $\frac{1}{2}||\mathbf{n}|| = \frac{1}{2}\sqrt{36+9+4} = \mathbf{7/2}$.

\subsection{Solutions to Q4}
\textbf{4a:} $\mathbf{n_1} = \langle 2,-1,1 \rangle, \mathbf{n_2} = \langle 1,1,2 \rangle$. $\cos\theta = \frac{|\mathbf{n_1} \cdot \mathbf{n_2}|}{||\mathbf{n_1}|| ||\mathbf{n_2}||} = \frac{2-1+2}{\sqrt{6}\sqrt{6}} = \frac{3}{6} = \frac{1}{2}$. \textbf{$\theta = \pi/3$}.\\
\textbf{4b:} $\mathbf{n_1} = \langle 1, \sqrt{3}, 2 \rangle, \mathbf{n_{xy}} = \langle 0,0,1 \rangle$. $\cos\theta = \frac{2}{\sqrt{1+3+4}(1)} = \frac{2}{\sqrt{8}} = \frac{1}{\sqrt{2}}$. \textbf{$\theta = \pi/4$}.\\
\textbf{4c:} $\mathbf{n_1} \cdot \mathbf{n_2} = 0 \implies (3)(4) + (k)(-1) + (-2)(k) = 12 - 3k = 0 \implies \mathbf{k = 4}$.\\
\textbf{4d:} $\mathbf{n_1} = \lambda \mathbf{n_2} \implies \langle a, 4, -6 \rangle = \lambda \langle 2, -b, 3 \rangle$. Looking at z: $-6 = 3\lambda \implies \lambda = -2$. $a = -4, 4 = 2b \implies \mathbf{a = -4, b = 2}$.

\subsection{Solutions to Q5}
\textbf{5a:} Area $= ||\mathbf{a} \times \mathbf{b}|| = || \langle (-2-12), (4 - (-4)), (6 - (-1)) \rangle || = || \langle -14, 8, 7 \rangle || = \sqrt{196+64+49} = \sqrt{309}$.\\
\textbf{5b:} $\vec{PQ} = \langle 2, 1, -1 \rangle, \vec{PR} = \langle 1, 3, 1 \rangle$. Area $= ||\vec{PQ} \times \vec{PR}|| = ||\langle 4, -3, 5 \rangle|| = \sqrt{16+9+25} = \sqrt{50} = \mathbf{5\sqrt{2}}$.\\
\textbf{5c:} Area $= \frac{1}{2} || \langle 1, -2, 3 \rangle \times \langle -1, 4, 2 \rangle || = \frac{1}{2} || \langle -16, -5, 2 \rangle || = \frac{1}{2}\sqrt{256+25+4} = \mathbf{\frac{\sqrt{285}}{2}}$.\\
\textbf{5d:} $\mathbf{u} \times \mathbf{v} = \langle 6, -3, c \rangle$. $||\mathbf{u} \times \mathbf{v}|| = \sqrt{36+9+c^2} = 7 \implies 45+c^2 = 49 \implies c^2 = 4 \implies \mathbf{c = \pm 2}$.

\subsection{Solutions to Q6}
\textbf{6a:} $\mathbf{v} = \langle 3, -4, 4 \rangle$. \textbf{$x = 2+3t, y = 4-4t, z = -1+4t$}.\\
\textbf{6b:} $\mathbf{v} = \langle -1, 2, 5 \rangle$. \textbf{$x = 1-t, y = 1+2t, z = 1+5t$}.\\
\textbf{6c:} $\mathbf{n} = \langle 4, -1, 7 \rangle$ becomes $\mathbf{v}$. \textbf{$x = -3+4t, y = 2-t, z = 6+7t$}.\\
\textbf{6d:} $\mathbf{v} = \mathbf{n_1} \times \mathbf{n_2} = \langle 1,1,1 \rangle \times \langle 1,-2,3 \rangle = \langle 3-(-2), 1-3, -2-1 \rangle = \mathbf{\langle 5, -2, -3 \rangle}$.

\subsection{Solutions to Q7}
\textbf{7a:} $\mathbf{f_x = 4x^3y^3 - 2e^{2x}}$, $\mathbf{f_y = 3x^4y^2 + 1/y}$.\\
\textbf{7b:} $\mathbf{f_x = 2x\sin(xy^3) + x^2y^3\cos(xy^3)}$, $\mathbf{f_y = 3x^3y^2\cos(xy^3)}$.\\
\textbf{7c:} $\mathbf{f_x = \frac{2(x+3y) - (2x-y)(1)}{(x+3y)^2} = \frac{7y}{(x+3y)^2}}$, $\mathbf{f_y = \frac{-1(x+3y) - (2x-y)(3)}{(x+3y)^2} = \frac{-7x}{(x+3y)^2}}$.\\
\textbf{7d:} $\mathbf{f_x = yz e^{xyz}\cos(x^2+z) - 2x e^{xyz}\sin(x^2+z)}$, $\mathbf{f_y = xz e^{xyz}\cos(x^2+z)}$, $\mathbf{f_z = xy e^{xyz}\cos(x^2+z) - e^{xyz}\sin(x^2+z)}$.

\subsection{Solutions to Q8}
\textbf{8a:} $f_x = 3x^2y^4 - 2y \implies \mathbf{f_{xx} = 6xy^4}$. $f_y = 4x^3y^3 - 2x \implies \mathbf{f_{yy} = 12x^3y^2}$. $\mathbf{f_{xy} = f_{yx} = 12x^2y^3 - 2}$.\\
\textbf{8b:} $f_x = y(e^{x^2y} \cdot 2xy) = 2xy^2e^{x^2y}$. $f_{xy} = \frac{\partial}{\partial y}(2xy^2e^{x^2y}) = 4xye^{x^2y} + 2xy^2e^{x^2y}(x^2) = \mathbf{4xye^{x^2y} + 2x^3y^2e^{x^2y}}$.\\
\textbf{8c:} $f_x = \frac{2x}{x^2+y^2} \implies f_{xy} = \frac{-2x(2y)}{(x^2+y^2)^2} = \frac{-4xy}{(x^2+y^2)^2}$. $f_y = \frac{2y}{x^2+y^2} \implies f_{yx} = \frac{-2y(2x)}{(x^2+y^2)^2} = \frac{-4xy}{(x^2+y^2)^2}$. They match.\\
\textbf{8d:} $u_x = e^x\sin y \implies u_{xx} = e^x\sin y$. $u_y = e^x\cos y \implies u_{yy} = -e^x\sin y$. $u_{xx} + u_{yy} = e^x\sin y - e^x\sin y = \mathbf{0}$.

\subsection{Solutions to Q9}
\textbf{9a:} $\nabla f = \langle 2xy, x^2 - 6y \rangle$. At $(1,2)$, $\nabla f = \langle 4, -11 \rangle$. $\mathbf{u} = \langle 3/5, 4/5 \rangle$. $D_{\mathbf{u}}f = 4(3/5) - 11(4/5) = -32/5 = \mathbf{-6.4}$.\\
\textbf{9b:} $\nabla f = \langle -ye^{-x}, e^{-x} \rangle$. At $(0,4)$, $\nabla f = \langle -4, 1 \rangle$. $\mathbf{u} = \langle \cos(2\pi/3), \sin(2\pi/3) \rangle = \langle -1/2, \sqrt{3}/2 \rangle$. $D_{\mathbf{u}}f = -4(-1/2) + 1(\sqrt{3}/2) = \mathbf{2 + \sqrt{3}/2}$.\\
\textbf{9c:} $\nabla f = \langle \frac{y}{2\sqrt{xy}}, \frac{x}{2\sqrt{xy}} \rangle$. At $(2,8)$, $\sqrt{16}=4$. $\nabla f = \langle 8/8, 2/8 \rangle = \langle 1, 1/4 \rangle$. Direction $\mathbf{v} = \langle -2, -8 \rangle \implies \mathbf{u} = \langle \frac{-2}{\sqrt{68}}, \frac{-8}{\sqrt{68}} \rangle$. $D_{\mathbf{u}}f = 1(\frac{-2}{\sqrt{68}}) + \frac{1}{4}(\frac{-8}{\sqrt{68}}) = \mathbf{\frac{-4}{\sqrt{68}}}$.\\
\textbf{9d:} $\nabla f = \langle 2xy, x^2+2yz, y^2 \rangle$. At $(1,2,-1)$, $\nabla f = \langle 4, 1-4, 4 \rangle = \langle 4, -3, 4 \rangle$. $\mathbf{u} = \langle 2/3, -1/3, 2/3 \rangle$. $D_{\mathbf{u}}f = 8/3 + 3/3 + 8/3 = \mathbf{19/3}$.

\subsection{Solutions to Q10}
\textbf{10a:} $\nabla f = \langle 3x^2, -3y^2 \rangle$. At $(2,-1)$, $\nabla f = \langle 12, -3 \rangle$. Max rate = $||\nabla f|| = \sqrt{144+9} = \mathbf{\sqrt{153}}$. Direction = $\mathbf{\langle 12, -3 \rangle}$.\\
\textbf{10b:} $\nabla f = \langle \frac{x}{\sqrt{x^2+2y}}, \frac{1}{\sqrt{x^2+2y}} \rangle$. At $(3,8)$, $\sqrt{9+16}=5$. $\nabla f = \langle 3/5, 1/5 \rangle$. Max rate = $\sqrt{9/25 + 1/25} = \sqrt{10/25} = \mathbf{\frac{\sqrt{10}}{5}}$.\\
\textbf{10c:} Min rate is in direction $-\nabla f$. $\nabla f = \langle e^{xy}+xy e^{xy}, x^2 e^{xy} \rangle$. At $(1,0)$, $\nabla f = \langle 1, 1 \rangle$. Direction = $\mathbf{\langle -1, -1 \rangle}$. Min rate = $\mathbf{-\sqrt{2}}$.\\
\textbf{10d:} $\nabla f = \langle 2x, 8y \rangle$. At $(2,1)$, $\nabla f = \langle 4, 8 \rangle$. We need $\mathbf{u} = \langle a, b \rangle$ such that $\langle 4, 8 \rangle \cdot \langle a, b \rangle = 0 \implies 4a+8b=0 \implies a=-2b$. Since $a^2+b^2=1 \implies 5b^2=1 \implies b = \pm 1/\sqrt{5}$. \textbf{$\mathbf{u} = \langle \frac{-2}{\sqrt{5}}, \frac{1}{\sqrt{5}} \rangle$ or $\langle \frac{2}{\sqrt{5}}, \frac{-1}{\sqrt{5}} \rangle$.}

\subsection{Solutions to Q11}
\textbf{11a:} $\int_0^2 \int_0^3 (x^2+y^2) dy dx = \int_0^2 [x^2y + y^3/3]_0^3 dx = \int_0^2 (3x^2 + 9) dx = [x^3 + 9x]_0^2 = 8 + 18 = \mathbf{26}$.\\
\textbf{11b:} Intersection: $x^2 = 4x \implies x=0, x=4$. $\int_0^4 \int_{x^2}^{4x} xy dy dx = \int_0^4 [\frac{1}{2}xy^2]_{x^2}^{4x} dx = \frac{1}{2} \int_0^4 (16x^3 - x^5) dx = \frac{1}{2}[4x^4 - x^6/6]_0^4 = \frac{1}{2}(1024 - 4096/6) = \frac{1}{2}(1024 - 682.66) = \mathbf{512/3}$.\\
\textbf{11c:} $\int_1^2 \int_0^{y^2} e^{x/y} dx dy = \int_1^2[y e^{x/y}]_0^{y^2} dy = \int_1^2 y(e^y - 1) dy$. Integration by parts: $[y e^y - e^y - y^2/2]_1^2 = (2e^2 - e^2 - 2) - (e - e - 1/2) = \mathbf{e^2 - 3/2}$.\\
\textbf{11d:} Region: $0 \le x \le 4, \sqrt{x} \le y \le 2 \implies 0 \le y \le 2, 0 \le x \le y^2$. Integral: $\int_0^2 \int_0^{y^2} \frac{1}{y^3+1} dx dy = \int_0^2 \frac{y^2}{y^3+1} dy = [\frac{1}{3} \ln(y^3+1)]_0^2 = \mathbf{\frac{1}{3} \ln(9)}$.

\subsection{Solutions to Q12}
\textbf{12a:} $\int_0^{\pi/2} \int_0^3 r^2 \cdot r dr d\theta = (\int_0^{\pi/2} d\theta) (\int_0^3 r^3 dr) = (\pi/2)(81/4) = \mathbf{\frac{81\pi}{8}}$.\\
\textbf{12b:} $\int_0^\pi \int_0^4 e^{r^2} r dr d\theta = \pi [\frac{1}{2}e^{r^2}]_0^4 = \mathbf{\frac{\pi}{2}(e^{16} - 1)}$.\\
\textbf{12c:} $\int_0^{2\pi} \int_1^2 \sin(r^2) r dr d\theta = 2\pi [-\frac{1}{2}\cos(r^2)]_1^2 = \mathbf{-\pi(\cos(4) - \cos(1))}$.\\
\textbf{12d:} $x^2+y^2 = 2x \implies r^2 = 2r\cos\theta \implies r = 2\cos\theta$. Area $= \int_{-\pi/2}^{\pi/2} \int_0^{2\cos\theta} 1 r dr d\theta = \mathbf{\pi}$.

\subsection{Solutions to Q13}
\textbf{13a:} Region is upper half cylinder. $\mathbf{\int_0^{\pi} \int_0^2 \int_0^{r^2} r^2\cos^2\theta \, r dz dr d\theta}$.\\
\textbf{13b:} Region is right half of a circle of radius 3 ($x$ goes from 0 to 3, $y$ from $- \dots$ to $+ \dots$). Note: bounded by $y$, so it's a circle in xy-plane. Wait, bounds are $dx dy$: $x \in [0, 3]$, $y \in[-\sqrt{9-x^2}, \sqrt{9-x^2}]$. This is the right half of circle. $-\pi/2 \le \theta \le \pi/2$. Integrand $xyz = (r\cos\theta)(r\sin\theta)z$. \textbf{$\int_{-\pi/2}^{\pi/2} \int_0^3 \int_0^r z r^3 \cos\theta \sin\theta \, dz dr d\theta$.}\\
\textbf{13c:} Region is a full sphere of radius 1. Cylindrical limits: $0 \le \theta \le 2\pi$, $0 \le r \le 1$, $-\sqrt{1-r^2} \le z \le \sqrt{1-r^2}$. Integrand $(x^2+y^2) = r^2$. \textbf{$\int_0^{2\pi} \int_0^1 \int_{-\sqrt{1-r^2}}^{\sqrt{1-r^2}} r^3 \, dz dr d\theta$.}\\
\textbf{13d:} Cylinder $r^2 = 2r\sin\theta \implies r = 2\sin\theta$. \textbf{$\int_0^{\pi} \int_0^{2\sin\theta} \int_0^{r^2} z r \, dz dr d\theta$.}

\subsection{Solutions to Q14}
\textbf{14a:} Region: first quadrant of circle radius 3. Integrand: $\frac{z}{r} r = z$. \textbf{$\int_0^{\pi/2} \int_0^3 \int_0^{r^2} z \, dz dr d\theta$.}\\
\textbf{14b:} Region: full circle radius 2. Integrand: $r \cdot r = r^2$. \textbf{$\int_0^{2\pi} \int_0^2 \int_{r^2}^4 r^2 \, dz dr d\theta$.}\\
\textbf{14c:} Region: first quadrant square $1 \times 1$? No, $x \in [0,1], y \in [0, \sqrt{1-x^2}]$. It's a quarter circle. \textbf{$\int_0^{\pi/2} \int_0^1 \int_0^1 e^{r^2} r \, dz dr d\theta$.}\\
\textbf{14d:} Region: right half circle $R=1$. Upper bound $z=1$, lower bound cone $z=r$. \textbf{$\int_{-\pi/2}^{\pi/2} \int_0^1 \int_r^1 z r \, dz dr d\theta$.}

\subsection{Solutions to Q15}
\textbf{15a:} $\int_0^{2\pi} \int_0^3 \int_0^{9-r^2} z r dz dr d\theta = 2\pi \int_0^3 [\frac{1}{2}z^2]_0^{9-r^2} r dr = \pi \int_0^3 (9-r^2)^2 r dr$. U-sub $u=9-r^2 \implies du=-2rdr$. $\pi \int_0^9 \frac{1}{2}u^2 du = \frac{\pi}{2}[\frac{u^3}{3}]_0^9 = \mathbf{\frac{243\pi}{2}}$.\\
\textbf{15b:} Intersection: $r=4$. $\int_0^{2\pi} \int_0^4 \int_r^4 r dz dr d\theta = 2\pi \int_0^4 r(4-r) dr = 2\pi [2r^2 - r^3/3]_0^4 = 2\pi(32 - 64/3) = \mathbf{64\pi/3}$.\\
\textbf{15c:} Intersection: $r^2 = 8-r^2 \implies 2r^2=8 \implies r=2$. $\int_0^{2\pi} \int_0^2 \int_{r^2}^{8-r^2} r^2 r dz dr d\theta = 2\pi \int_0^2 r^3(8-2r^2) dr = 2\pi[2r^4 - r^6/3]_0^2 = 2\pi(32 - 64/3) = \mathbf{64\pi/3}$.\\
\textbf{15d:} Mass $M = \iiint r dz dr d\theta = 64\pi/3$. Moment $M_{xy} = \iiint z r dz dr d\theta = \int_0^{2\pi} \int_0^4 \int_r^4 z r dz dr d\theta = 2\pi \int_0^4 \frac{r}{2}(16-r^2) dr = \pi [8r^2 - r^4/4]_0^4 = \pi(128 - 64) = 64\pi$. $\bar{z} = M_{xy} / M = \mathbf{3}$.

\subsection{Solutions to Q16}
\textbf{16a:} Region: Full circle $R=2$, $z$ from 0 to upper hemisphere. This is a top hemisphere of radius 2. Integrand: $(\rho\cos\phi)^2 \rho \cdot \rho^2\sin\phi = \rho^5\cos^2\phi\sin\phi$. \textbf{$\int_0^{2\pi} \int_0^{\pi/2} \int_0^2 \rho^5 \cos^2\phi \sin\phi \, d\rho d\phi d\theta$.}\\
\textbf{16b:} First octant sphere $R=3$. Integrand: $x \rho^2\sin\phi = \rho^3\sin^2\phi\cos\theta$. \textbf{$\int_0^{\pi/2} \int_0^{\pi/2} \int_0^3 \rho^3\sin^2\phi\cos\theta \, d\rho d\phi d\theta$.}\\
\textbf{16c:} $z$ from $r$ (cone $\phi=\pi/4$) to sphere $R=\sqrt{2}$. Full rotation. Integrand: $\frac{1}{\rho^2} \rho^2\sin\phi = \sin\phi$. \textbf{$\int_0^{2\pi} \int_0^{\pi/4} \int_0^{\sqrt{2}} \sin\phi \, d\rho d\phi d\theta$.}\\
\textbf{16d:} Below $z=1 \implies \rho\cos\phi = 1 \implies \rho = \sec\phi$. Intersection of $z=1$ and sphere is $\rho=2, \cos\phi = 1/2 \implies \phi = \pi/3$. Bounds: $0 \le \theta \le 2\pi$, $0 \le \phi \le \pi/3$, $\sec\phi \le \rho \le 2$. \textbf{$\int_0^{2\pi} \int_0^{\pi/3} \int_{\sec\phi}^2 (\rho\cos\phi) \rho^2\sin\phi \, d\rho d\phi d\theta$.}

\subsection{Solutions to Q17}
\textbf{17a:} Inside sphere $\rho=5$, inside cone $\phi=\pi/4$. Volume $= \int_0^{2\pi} \int_0^{\pi/4} \int_0^5 \rho^2\sin\phi \, d\rho d\phi d\theta$.\\
\textbf{17b:} First octant $\implies 0 \le \theta \le \pi/2$. Spheres limit $\rho \in [1,3]$. Cones limit $\phi \in[\pi/6, \pi/3]$. \textbf{$\int_0^{\pi/2} \int_{\pi/6}^{\pi/3} \int_1^3 f(\rho,\theta,\phi) \rho^2\sin\phi \, d\rho d\phi d\theta$.}\\
\textbf{17c:} "Bounded below by xy-plane ($\phi=\pi/2$), sides by sphere $\rho=6$, above by cone $\phi=\pi/6$". Bounds: $\pi/6 \le \phi \le \pi/2, 0 \le \rho \le 6, 0 \le \theta \le 2\pi$. $V = \int_0^{2\pi} \int_{\pi/6}^{\pi/2} \int_0^6 \rho^2\sin\phi \, d\rho d\phi d\theta = 2\pi [-\cos\phi]_{\pi/6}^{\pi/2} [\rho^3/3]_0^6 = 2\pi (0 - (-\sqrt{3}/2)) (72) = \mathbf{72\pi\sqrt{3}}$.\\
\textbf{17d:} $\rho^2 = 2\rho\cos\phi \implies \mathbf{\rho = 2\cos\phi}$. Because $\rho \ge 0$, $\cos\phi \ge 0 \implies \mathbf{0 \le \phi \le \pi/2}$. $0 \le \theta \le 2\pi$.

\subsection{Solutions to Q18}
\textbf{18a:} $P_y = 3, Q_x = 5$. $\iint_D (5-3) dA = 2(\text{Area}) = 2(2) = \mathbf{4}$.\\
\textbf{18b:} $P_y = 2y, Q_x = 3y$. Clockwise $\implies$ Negative Green's. $-\iint_D (3y-2y) dA = -\iint_D y dA$. Since region is symmetric circle over origin, $\iint y dA = 0$. \textbf{Answer: 0}.\\
\textbf{18c:} $P_y = -2, Q_x = 4$. $\iint_D (4 - (-2)) dA = 6(\text{Area})$. Area of triangle $= 1/2(4)(3) = 6$. \textbf{Answer: 36}.\\
\textbf{18d:} $A = \frac{1}{2} \int_0^{2\pi} (\cos^3 t)(3\sin^2 t \cos t) - (\sin^3 t)(-3\cos^2 t \sin t) dt = \frac{1}{2} \int_0^{2\pi} 3\sin^2 t \cos^2 t (\cos^2 t + \sin^2 t) dt = \frac{3}{2} \int_0^{2\pi} \frac{1}{4}\sin^2(2t) dt = \frac{3}{8}[\frac{1}{2}t - \frac{1}{8}\sin(4t)]_0^{2\pi} = \mathbf{\frac{3\pi}{8}}$.

\newpage
% ---------------------------------------------------------
% ROADMAPS & MENTAL MODELS
% ---------------------------------------------------------
\section{Roadmaps \& Mental Checklists for Q1-Q18}

\subsection{Roadmap: 2D & 3D Double Integrals (Q11-Q12)}
\textbf{The Mental Trigger:} You are asked to find the volume under a surface $z = f(x,y)$, or evaluate an integral $\iint_D f(x,y)dA$.
\begin{itemize}
    \item \textbf{Check 1 (Shape of Domain $D$):} Look at the bounds given.
    \begin{itemize}
        \item \textbf{IF} the bounds are rectangular ($x \in [a,b], y \in [c,d]$): Use standard Cartesian $dx dy$.
        \item \textbf{IF} the bounds involve circles, semi-circles, or $x^2+y^2$: IMMEDIATELY switch to Polar.
    \end{itemize}
    \item \textbf{Check 2 (Polar Conversion Checklist):}
    \begin{enumerate}
        \item Replace all $x^2+y^2$ with $r^2$.
        \item Replace $x$ with $r\cos\theta$ and $y$ with $r\sin\theta$.
        \item \textbf{CRITICAL:} Replace $dA$ with $r dr d\theta$. Do not forget the Jacobian $r$!
        \item Graph the 2D region $D$ to find bounds.
        \begin{itemize}
            \item Full circle: $\theta \in [0, 2\pi]$
            \item Upper half: $\theta \in [0, \pi]$
            \item First quadrant: $\theta \in[0, \pi/2]$
            \item Shifted circle (e.g., $x^2+y^2=2x \implies r=2\cos\theta$): $\theta \in[-\pi/2, \pi/2]$, $r \in [0, 2\cos\theta]$.
        \end{itemize}
    \end{enumerate}
\end{itemize}

\subsection{Roadmap: Triple Integrals in Cylindrical Coordinates (Q13-Q15)}
\textbf{The Mental Trigger:} You are dealing with 3D volumes or masses, and the solid involves cylinders ($x^2+y^2=R^2$), paraboloids ($z=x^2+y^2$), or cones ($z=\sqrt{x^2+y^2}$).
\begin{itemize}
    \item \textbf{Step 1: Set the $z$-bounds.} 
    \begin{itemize}
        \item Identify the "floor" (lower surface) and "ceiling" (upper surface).
        \item Convert their Cartesian equations to cylindrical (replace $x^2+y^2$ with $r^2$).
        \item Example: Between $z=0$ and $z=4-x^2-y^2 \implies$ Bounds are $0 \le z \le 4-r^2$.
    \end{itemize}
    \item \textbf{Step 2: Find the "Shadow" (xy-plane bounds).}
    \begin{itemize}
        \item Set the floor and ceiling equations equal to each other to find their intersection.
        \item This intersection creates a 2D region $D$ in the xy-plane. 
        \item Determine the $r$ and $\theta$ bounds for this 2D region exactly like a Polar Double Integral.
    \end{itemize}
    \item \textbf{Step 3: Setup the Integral.}
    \begin{itemize}
        \item $\iiint f(r,\theta,z) \cdot \mathbf{r} \, dz \, dr \, d\theta$.
        \item Always evaluate $dz$ first, then $dr$, then $d\theta$.
    \end{itemize}
\end{itemize}

\subsection{Roadmap: Triple Integrals in Spherical Coordinates (Q16-Q17)}
\textbf{The Mental Trigger:} The solid involves spheres ($x^2+y^2+z^2 = R^2$) or pure cones ($z=\sqrt{x^2+y^2}$). The integrand often contains $(x^2+y^2+z^2)^{n}$.
\begin{itemize}
    \item \textbf{Step 1: The Jacobian.} Immediately write $dV = \mathbf{\rho^2 \sin\phi \, d\rho \, d\phi \, d\theta}$.
    \item \textbf{Step 2: Determine $\rho$ (Radius from origin).}
    \begin{itemize}
        \item Solid sphere of radius $R \implies \rho \in [0, R]$.
        \item Hollow sphere (between radius 2 and 4) $\implies \rho \in [2, 4]$.
        \item Flat plane top ($z=a$) $\implies \rho\cos\phi = a \implies \rho = a\sec\phi$.
    \end{itemize}
    \item \textbf{Step 3: Determine $\phi$ (Angle dropping from positive z-axis).}
    \begin{itemize}
        \item Positive z-axis $\implies \phi = 0$.
        \item xy-plane $\implies \phi = \pi/2$.
        \item Upper half-space (above xy-plane) $\implies \phi \in [0, \pi/2]$.
        \item A cone $z = \sqrt{x^2+y^2} \implies \phi = \pi/4$.
        \item "Inside" the cone means $0 \le \phi \le \pi/4$. "Outside" the cone means $\pi/4 \le \phi \le \pi/2$.
    \end{itemize}
    \item \textbf{Step 4: Determine $\theta$ (Standard xy-plane rotation).}
    \begin{itemize}
        \item Full 3D object $\implies \theta \in [0, 2\pi]$.
        \item First octant ($x>0, y>0$) $\implies \theta \in [0, \pi/2]$.
    \end{itemize}
\end{itemize}

\subsection{Roadmap: Green's Theorem (Q18)}
\textbf{The Mental Trigger:} You are asked to evaluate a line integral $\oint_C P dx + Q dy$ over a \textbf{CLOSED LOOP} $C$ in 2D.
\begin{itemize}
    \item \textbf{Step 1: Check Orientation.}
    \begin{itemize}
        \item \textbf{IF} Counterclockwise (CCW): Standard application.
        \item \textbf{IF} Clockwise (CW): You MUST multiply your final double integral by $-1$.
    \end{itemize}
    \item \textbf{Step 2: Calculate the Curl.}
    \begin{itemize}
        \item Find $\frac{\partial Q}{\partial x} - \frac{\partial P}{\partial y}$.
        \item \textbf{TRAP:} Do not mix them up. It is NOT $P_x - Q_y$. 
    \end{itemize}
    \item \textbf{Step 3: Set up Double Integral.}
    \begin{itemize}
        \item $\iint_D (\frac{\partial Q}{\partial x} - \frac{\partial P}{\partial y}) dA$.
        \item If the curl is a constant $K$, the answer is simply $K \times \text{Area}(D)$.
        \item If the region $D$ is a circle, immediately convert the double integral to polar coordinates.
    \end{itemize}
\end{itemize}

\end{document}